% Generated draft: requires mathematical and source review.
\documentclass[11pt]{article}
\usepackage[margin=25mm]{geometry}
\usepackage{fontspec}
\setmainfont{Latin Modern Roman}
\usepackage{amsmath,amssymb,mathtools,unicode-math}
\setmathfont{Latin Modern Math}
\usepackage{enumitem,needspace,xurl,hyperref}
\setlength{\parindent}{0pt}
\setlength{\parskip}{5pt}
\setlength{\emergencystretch}{3em}
\newcommand{\N}{\mathbb{N}}
\newcommand{\Z}{\mathbb{Z}}
\newcommand{\Q}{\mathbb{Q}}
\newcommand{\R}{\mathbb{R}}
\newcommand{\C}{\mathbb{C}}
\newcommand{\F}{\mathbb{F}}
\newcommand{\A}{\mathbb{A}}
\newcommand{\PP}{\mathbb{P}}
\newcommand{\E}{\mathbb{E}}
\newcommand{\eps}{\varepsilon}
\newcommand{\dd}{\,\mathrm{d}}
\begin{document}
\section*{UnsolvedMath ID: 10600031}
\textbf{Source problem code:} catalogue entry
\subsection{Definitions and mathematical statement}
\emph{Research problem.} Make a systematic study of Vassiliev invariants for virtual knots and links. The source states that work on this program is ongoing; it does not give a single conjectural identity or a finite list of classification targets.

For a chosen class \(\mathcal K\) of virtual knots or links and a coefficient group \(A\), a finite-type theory consists of invariant functions \(v:\mathcal K\to A\) together with a specified singular-diagram extension or difference operation. In the positive-minus-negative crossing convention, one sets \(v(D_\times)=v(D_+)-v(D_-)\) and asks for vanishing after sufficiently many singularizations. Virtual-knot theory also admits other difference conventions, so this is an explanatory example, not a choice imposed by the supplied statement. A systematic formulation must first fix the virtual finite-type convention, its degree filtration, coefficients, and the knot or link equivalence being used, and then specify which structures or classification results are to be determined. The excerpt supplies none of those additional choices. It includes both knots and links and does not restrict the number of link components.
\subsection{Short English statement}
Develop a systematic theory of virtual-knot and virtual-link Vassiliev invariants. The source presents an ongoing research program, with no fixed coefficient system, finite-type convention, or complete classification target.
\subsection{Sources}
\begin{itemize}
\item Local UnsolvedMath registry, permanent ID 10600031, source code catalogue entry, ``Virtual-knot problem 31 --- Make a systematic study of Vassiliev invariants for virtual knots and links.'' No individual proposer is supplied.
\item \emph{Unsolved Problems in Virtual Knot Theory and Combinatorial Knot Theory}, revised survey, problem-list item 31. Supplied URL: \url{https://arxiv.org/abs/1409.2823}; recorded access 2026-07-29; extraction source-tex. The dedicated section contains 60 top-level enumerated items with compound items retained intact.
\item The research summary lists ``On Vassiliev Invariants of Virtual Knots'' (2022), \texttt{arXiv:2208.xxxx}; ``A Lattice of Finite-Type Invariants of Virtual Knots'' (2013), \texttt{arXiv:1303.xxxx}; ``Some Dimensions of Spaces of Finite Type Invariants of Virtual Knots'' (2009), \texttt{arXiv:0909.xxxx}; ``Vassiliev Invariants from Parity Mappings'' (2012), \texttt{arXiv:1203.xxxx}; and ``Parity and Exotic Combinatorial Formulae for Finite-Type Invariants of Virtual Knots'' (2010), \texttt{arXiv:1002.xxxx}. These are incomplete placeholder identifiers, not verified usable references.
\item Supplied status partially\_solved and PARTIAL-PROGRESS; the summary says systematic work is ongoing and no completion is documented. Current primary-literature status and redistribution terms of the publicly accessible author/publisher copy remain NEEDS\_REVIEW.
\end{itemize}
\end{document}
